%% =====================================================================
%%  T-05-03  The Black Swan
%%  -------------------------------------------------------------------
%%  One idea per file. Core is mandatory. Slots in canonical order.
%%  Max 700 words. No layout commands. No filler. New source material
%%  for this entry goes in sources/<BOOK>/T-05-03.tex -- never by
%%  rewriting this file (docs/RULES.md R0.1).
%% =====================================================================
%% @kind: topic
%% @id: T-05-03
%% @title: The Black Swan
%% @subtitle: The fact you did not know you did not know, and what it re-prices
%% @chapter: c05
%% @order: 30
%% @status: stable
%% @sources: B5
%% @updated: 2026-09-19
%% @allow-rewrite: B5 ch. 8-10 ingest 2026-09-19 -- committed scaffold replaced by the written entry (planned -> stable lifecycle)

\topic{T-05-03}{The Black Swan}{The fact you did not know you did not know, and what it re-prices}

\begin{core}
Every negotiation runs on a picture of the counterpart's world, and the load-bearing fact of that world is often one you have not even failed to learn --- an unknown unknown. When it surfaces, it does not amend the deal; it re-prices everything. Hunting for that fact is worth more than optimising inside the assumptions you already hold.
\end{core}
\begin{mechanism}
Assumptions operate as if they were data. The known knowns --- the case file, the stated positions --- get all the attention, and over-reliance shackles a negotiator to a picture that cannot update. The unknown unknown stay invisible precisely because nothing prompts their question. They matter disproportionately because deals are structured around the counterpart's real constraint --- the hidden motive, the unspoken deadline, the third party --- and one surfaced fact re-prices every number built on it. Getting it out depends on two documented levers. Similarity unlocks disclosure --- people trust an in-group reader --- and disclosure lives in face time; it trickles in thin channels. The act of surfacing is mostly receptive --- curiosity about their world beyond the deal, and the discipline of chasing surprises instead of smoothing them.
\end{mechanism}
\begin{conditions}
Pays most where stakes are high, the counterpart is opaque, and talks are deadlocked on positions --- there a single hidden fact can break a year of stalemate. Pays least in transparent, repeated, small transactions, where the picture is cheap to verify and the swans have already flown. A counterpart who knows you hunt swans can feed you one; a fact that arrives too conveniently is bait until verified (T-23-03).
\end{conditions}
\begin{application}
  \item Before contact, write your assumptions as assumptions: what am I certain of that I have never checked?
  \item Ask about their world beyond the deal --- how they came to this, what else it touches, who else it must survive.
  \item Build real similarity where it exists --- shared background, taste, enemy --- and never fake it; counterfeit kinship is worse than none (T-23-02).
  \item Prefer richest-channel contact: in person over voice, voice over text.
  \item When something surprises you, chase it --- a surprise is the sound of an assumption breaking.
\end{application}
\begin{feedback}
  \item Landing: the counterpart corrects you gently (\emph{actually, the reason we...}) --- an assumption just died in front of them.
  \item Landing: warmth rises after a genuine commonality surfaces, and specifics start arriving unprompted.
  \item Failing: flat repetition of rehearsed positions --- no disclosure channel is open yet.
  \item Failing: alarm at your curiosity; the question has touched the thing they are protecting.
\end{feedback}
\begin{failure}
Faked or batched curiosity degrades the hunt into interrogation and closes the channel it depends on. One juicy fact can also over-fit: re-pricing everything on a single reveal is how a planted fact manipulates you, and the bigger the swan, the more it needs verification first. The subtlest cost is self-inflicted --- a hunter in love with the hunt stops hearing the ordinary facts that were never hidden.
\end{failure}
\begin{countermeasures}
  \item Audit the other direction: what do they know about you that you have not surfaced? Your own unknown unknowns are leaking somewhere.
  \item Treat every volunteered \emph{one thing you should know} as evidence, not as a gift --- verify it like any other claim.
  \item Notice installed kinship. The in-group feeling can be manufactured on purpose, and its purpose is your disclosure (T-22-02).
  \item When one fact flips your whole read, sleep on it. Re-pricing is allowed; instant re-pricing is how bait works.
\end{countermeasures}

\sources{B5}
\seesources
\seealso{T-05-01, T-05-02, T-07-01}
